%% ch02.tex --- Linear Regression with One Variable
%% Numbers are computed by verify/ch02_cost_function.py.

\chapter{Linear Regression with One Variable}
\label{ch:ch02}

\epigraph{All models are wrong, but some are useful.}{---\ George E.\ P.\ Box}

\begin{chapterlearning}
After working through this chapter you will be able to:
\begin{enumerate}
  \item Write the univariate linear regression model and say what each of its
        two parameters does to the fitted line.
  \item Define the squared-error cost function, and justify each of its three
        design choices: the square, the $1/m$, and the $\tfrac12$.
  \item Compute $J(w,b)$ by hand for a given dataset and a given pair of
        parameters.
  \item Explain why $J$ is a parabola in one parameter and a bowl in two, and
        why that guarantees a single minimum.
  \item Read a contour plot of $J$: identify the minimum, compare two parameter
        pairs, and recognise what elongated contours indicate.
  \item Derive the closed-form solution for the best-fitting line, and apply it.
  \item Explain why a low cost does not by itself mean a good model.
\end{enumerate}
\end{chapterlearning}

\Cref{ch:ch01} said that every supervised learner has four parts: data, a
hypothesis space, a cost, and an optimizer. This chapter examines the third of
them in the simplest setting that exists --- fitting a straight line ---
because the squared-error cost introduced here reappears, essentially unchanged,
in \cref{ch:ch04,ch:ch05}, is modified only in its inner expression in
\cref{ch:ch06}, and is what neural networks in \cref{ch:ch09} are still
minimising. Understand it once, thoroughly, and much of the rest of the book is
bookkeeping.

\section{The Model}
\label{sec:the-model}

\begin{defbox}{Univariate linear regression}
Given a single input feature $x$, the \term{univariate linear regression} model
is
\[
  \fwb(x) \;=\; wx + b ,
\]
with parameters $w\in\R$, the \term{weight} (or slope), and $b\in\R$, the
\term{bias} (or intercept). The hypothesis space is
$\Hspace = \{\,x \mapsto wx+b \;:\; w,b\in\R\,\}$: one line for each pair
$(w,b)$.
\end{defbox}

The two parameters do different jobs, and it is worth fixing the picture before
any algebra. Changing $w$ tilts the line about its intercept; changing $b$
slides it vertically without tilting it. \Cref{fig:wb-effect} shows both.

\begin{figure}[htbp]
  \centering
  \begin{minipage}[t]{0.47\linewidth}
    \centering
    \begin{tikzpicture}
      \begin{axis}[mlaxis, width=\linewidth, height=4.2cm,
        xlabel={$x$}, ylabel={$y$}, xmin=0, xmax=4, ymin=-1, ymax=7,
        legend pos=north west, legend style={font=\tiny}]
        \addplot[MLBlue,thick,domain=0:4]{0.5*x+1}; \addlegendentry{$w=0.5$}
        \addplot[MLTeal,thick,domain=0:4]{1.5*x+1}; \addlegendentry{$w=1.5$}
        \addplot[MLRose,thick,domain=0:4]{2.5*x+1}; \addlegendentry{$w=2.5$}
      \end{axis}
    \end{tikzpicture}
    {\scriptsize Varying $w$ with $b=1$ fixed: the line pivots.\par}
  \end{minipage}\hfill
  \begin{minipage}[t]{0.47\linewidth}
    \centering
    \begin{tikzpicture}
      \begin{axis}[mlaxis, width=\linewidth, height=4.2cm,
        xlabel={$x$}, ylabel={$y$}, xmin=0, xmax=4, ymin=-1, ymax=7,
        legend pos=north west, legend style={font=\tiny}]
        \addplot[MLBlue,thick,domain=0:4]{1.2*x-0.5}; \addlegendentry{$b=-0.5$}
        \addplot[MLTeal,thick,domain=0:4]{1.2*x+1};   \addlegendentry{$b=1$}
        \addplot[MLRose,thick,domain=0:4]{1.2*x+2.5}; \addlegendentry{$b=2.5$}
      \end{axis}
    \end{tikzpicture}
    {\scriptsize Varying $b$ with $w=1.2$ fixed: the line shifts.\par}
  \end{minipage}
  \caption{The two parameters of the model, seen separately. Every straight line
  in the plane is one point $(w,b)$, and vice versa --- a correspondence we will
  exploit heavily from \cref{sec:cost-two-params} onward.}
  \label{fig:wb-effect}
\end{figure}

That last sentence in the caption is the idea that makes this chapter work.
There are two spaces in play, and we will move between them constantly:

\begin{itemize}
  \item \textbf{Data space}, with $x$ horizontal and $y$ vertical. A model is a
        \emph{line}; a training example is a \emph{point}.
  \item \textbf{Parameter space}, with $w$ horizontal and $b$ vertical. A model
        is a \emph{point}; the training set is fixed and invisible, having been
        absorbed into the cost.
\end{itemize}

Confusing the two is the single most common source of bewilderment in this
material. When you see a plot, ask first what its axes are.

\section{Notation for the Training Set}
\label{sec:notation}

We are given $m$ examples $(\xii{1},\yii{1}),\dots,(\xii{m},\yii{m})$. Throughout
this chapter the running example is four points, small enough to compute with by
hand:

\begin{center}
\small
\renewcommand{\arraystretch}{1.2}
\begin{tabular}{@{}lrrrr@{}}
\toprule
$i$      & $1$ & $2$ & $3$ & $4$\\
\midrule
$\xii{i}$ & $1$ & $2$ & $3$ & $4$\\
$\yii{i}$ & $1$ & $3$ & $4$ & $6$\\
\bottomrule
\end{tabular}
\end{center}

\noindent so $m = 4$, $\xii{2} = 2$ and $\yii{2} = 3$. Call this the
\term{four-point set}; it is dataset~T4 in \cref{ch:appD}. Its points are not
collinear --- no straight line passes through all four --- which is what makes
the fitting problem non-trivial and, as it happens, what makes every real
regression problem non-trivial.

\section{Measuring Fit: The Cost Function}
\label{sec:cost}

Fix a candidate line, say $\fwb(x) = x + 1$, that is $(w,b) = (1,1)$. For each
training example it makes a prediction $\yhat^{(i)} = \fwb(\xii{i})$, and each
prediction misses its target by some amount.

\begin{defbox}{Residual}
The \term{residual} of example $i$ under the model $\fwb$ is
\[
  r^{(i)} \;=\; \yhat^{(i)} - \yii{i} \;=\; \fwb(\xii{i}) - \yii{i} .
\]
It is positive when the model over-predicts and negative when it
under-predicts.
\end{defbox}

\begin{figure}[htbp]
  \centering
  \begin{tikzpicture}
    \begin{axis}[mlaxis, width=0.62\linewidth, height=5.0cm,
      xlabel={$x$}, ylabel={$y$}, xmin=0, xmax=5, ymin=0, ymax=7.5,
      legend pos=north west, legend style={font=\tiny}]
      \addplot[MLBlue,thick,domain=0.3:4.6]{1*x+1};
      \addlegendentry{$\fwb(x)=x+1$}
      \addplot[only marks,mark=*,mark size=1.8pt,MLTeal]
        coordinates {(1,1)(2,3)(3,4)(4,6)};
      \addlegendentry{training examples}
      \addplot[MLRose,thick,dashed,forget plot] coordinates {(1,1)(1,2)};
      \addplot[MLRose,thick,dashed,forget plot] coordinates {(2,3)(2,3)};
      \addplot[MLRose,thick,dashed,forget plot] coordinates {(3,4)(3,4)};
      \addplot[MLRose,thick,dashed,forget plot] coordinates {(4,6)(4,5)};
      \node[font=\tiny,MLRose,anchor=west] at (axis cs:1.08,1.5)  {$r=+1$};
      \node[font=\tiny,MLRose,anchor=west] at (axis cs:2.08,3.0)  {$r=0$};
      \node[font=\tiny,MLRose,anchor=west] at (axis cs:3.08,4.0)  {$r=0$};
      \node[font=\tiny,MLRose,anchor=west] at (axis cs:4.08,5.5)  {$r=-1$};
    \end{axis}
  \end{tikzpicture}
  \caption{Residuals of the line $(w,b)=(1,1)$ on the four-point set: the
  vertical gaps between prediction and observation. The cost aggregates these
  four numbers into one.}
  \label{fig:residuals}
\end{figure}

We now need to collapse $m$ residuals into a single number, so that two
candidate lines can be compared. The standard choice is their mean square.

\begin{defbox}{Squared-error cost}
The \term{squared-error cost} of the model $\fwb$ on a training set of $m$
examples is
\[
  J(w,b) \;=\; \frac{1}{2m}\sum_{i=1}^{m}
              \bigl(\fwb(\xii{i}) - \yii{i}\bigr)^{2}
        \;=\; \frac{1}{2m}\sum_{i=1}^{m}\bigl(r^{(i)}\bigr)^{2} .
\]
It is a function of the parameters $(w,b)$; the data is held fixed.
\end{defbox}

\begin{derivation}{The cost as a squared length}
Collect the residuals into a vector $\mathbf{r}\in\Real{m}$ with components
$r^{(i)}$. Then
\[
  \sum_{i=1}^{m}\bigl(r^{(i)}\bigr)^{2}
  \;=\; \mathbf{r}\cdot\mathbf{r}
  \;=\; \|\mathbf{r}\|^{2},
  \qquad\text{so}\qquad
  J(w,b) \;=\; \frac{1}{2m}\,\|\mathbf{r}\|^{2}
         \;=\; \frac{1}{2m}\,\bigl\|\hat{\mathbf{y}} - \vy\bigr\|^{2}.
\]
The cost is therefore the squared Euclidean distance between the vector of
predictions and the vector of observations, scaled. Minimising $J$ means
bringing the prediction vector as close as possible to the observation vector in
$\Real{m}$ --- a reading we return to in \cref{sec:projection}, where it turns
out to explain the whole method geometrically.
\end{derivation}

Three features of this definition are chosen rather than forced, and each
deserves a justification.

\subsection{Why square the residuals}

The obvious first attempt, summing the residuals themselves, fails immediately:
positive and negative errors cancel. In \cref{fig:residuals} the residuals are
$+1, 0, 0, -1$, summing to zero, which would rate that mediocre line as perfect.
Any usable measure must treat over- and under-prediction as equally bad, so it
must be built from a symmetric, non-negative function of each residual.

That leaves two natural candidates, $|r|$ and $r^{2}$, and squares win for three
reasons:

\begin{enumerate}
  \item \textbf{Differentiability.} $|r|$ has no derivative at $r = 0$, exactly
        the point the optimizer is heading for. $r^{2}$ is smooth everywhere,
        and \cref{ch:ch03}'s method needs derivatives.
  \item \textbf{A closed-form solution exists.} Squaring makes $J$ a quadratic
        in the parameters, which we can minimise exactly with algebra
        (\cref{sec:closed-form}). Absolute values do not permit this.
  \item \textbf{Large errors dominate.} An error of $4$ contributes sixteen
        times as much as an error of $1$, so the fitted line works hard to avoid
        being badly wrong anywhere. Whether that is desirable is a modelling
        decision, not a mathematical fact --- see \cref{sec:low-cost}.
\end{enumerate}

There is also a statistical justification, which \cref{ch:ch06} makes precise:
if the observations differ from the true line by independent Gaussian noise,
the line minimising squared error is the one of maximum likelihood. Squaring is
not arbitrary; it is the correct thing to do under a specific and often
reasonable assumption about how the data was generated.

\subsection{Why divide by \texorpdfstring{$m$}{m}}

Without the $1/m$, the cost would grow with the size of the training set, and a
model fitted to $1000$ examples would appear a thousand times worse than the
same model fitted to one. Averaging makes the number comparable across datasets
of different sizes --- and, more immediately, means that adding more of the same
data does not change the answer.

\begin{workedexample}{The $1/m$ earns its place}
Take the four-point set, and a second dataset consisting of the same four points
each written down twice, so $m = 8$. Any sensible measure of fit should be
unchanged: nothing new has been observed.

\begin{center}
\small
\begin{tabular}{@{}lrrrr@{}}
\toprule
& \multicolumn{2}{c}{$\sum_i (r^{(i)})^2$} & \multicolumn{2}{c}{$J$}\\
\cmidrule(lr){2-3}\cmidrule(lr){4-5}
$(w,b)$ & $m=4$ & $m=8$ & $m=4$ & $m=8$\\
\midrule
$(1.6,-0.5)$ & $0.200$ & $0.400$ & $0.0250$ & $0.0250$\\
$(1.0,\ \ 1.0)$ & $2.000$ & $4.000$ & $0.2500$ & $0.2500$\\
\bottomrule
\end{tabular}
\end{center}

The raw sum doubles; the cost does not. Note also that the \emph{ranking} of the
two lines is unaffected either way --- which is why the $1/m$ is a convenience
rather than a necessity, and why you will occasionally see it omitted.
\end{workedexample}

\subsection{Why the one half}

The $\tfrac12$ has no effect whatsoever on which $(w,b)$ minimises $J$: scaling
a function by a positive constant moves its graph but not the location of its
lowest point. It is there purely so that the $2$ produced by differentiating the
square cancels it. \Cref{ch:ch03} derives
\[
  \frac{\partial J}{\partial w}
    = \frac{1}{m}\sum_{i=1}^{m}\bigl(\fwb(\xii{i})-\yii{i}\bigr)\xii{i},
\]
with no stray factor of two, and you will be grateful for the tidiness a hundred
gradient computations from now.

\begin{keyinsight}{What $J$ is a function of}
$J$ is a function of the \emph{parameters}, not of the data. When we write
$J(w,b)$ and plot it, the training set has been fixed and folded into the
function's definition. Every point of that plot answers one question: if the
line were this one, how badly would it fit the data we have?
\end{keyinsight}

\section{The Cost as a Function of One Parameter}
\label{sec:cost-one-param}

Two parameters are one too many for a first picture, so temporarily fix $b = 0$
and consider the simpler model $f_w(x) = wx$. Take three training points lying
exactly on a line through the origin: $(1,1)$, $(2,2)$, $(3,3)$, so that we know
in advance that $w = 1$ is perfect.

\begin{derivation}{$J(w)$ is a parabola}
With $b=0$, $m=3$ and $\yii{i} = \xii{i}$,
\[
  \begin{aligned}
    J(w) &= \frac{1}{2\cdot 3}\sum_{i=1}^{3}\bigl(w\xii{i}-\xii{i}\bigr)^{2}
          = \frac{1}{6}(w-1)^{2}\sum_{i=1}^{3}\bigl(\xii{i}\bigr)^{2}\\[2pt]
         &= \frac{1+4+9}{6}(w-1)^{2}
          = \frac{14}{6}(w-1)^{2}.
  \end{aligned}
\]
This is a parabola in $w$: it opens upward, because the coefficient
$\tfrac{14}{6}$ is positive, and it touches zero at $w=1$. Upward-opening
parabolas have exactly one minimum and no other stationary point, so for this
data the search for the best $w$ has a unique answer and no way to get stuck.
\end{derivation}

\begin{figure}[htbp]
  \centering
  \begin{minipage}[t]{0.47\linewidth}
    \centering
    \begin{tikzpicture}
      \begin{axis}[mlaxis, width=\linewidth, height=4.4cm,
        xlabel={$x$}, ylabel={$y$}, xmin=0, xmax=3.6, ymin=0, ymax=5,
        legend pos=north west, legend style={font=\tiny}]
        \addplot[only marks,mark=*,mark size=1.8pt,MLTeal]
          coordinates {(1,1)(2,2)(3,3)};
        \addlegendentry{data}
        \addplot[MLRose,thick,domain=0:3.4]{0.5*x};  \addlegendentry{$w=0.5$}
        \addplot[MLBlue,very thick,domain=0:3.4]{1*x}; \addlegendentry{$w=1$}
        \addplot[MLAmber,thick,domain=0:3.4]{1.5*x}; \addlegendentry{$w=1.5$}
      \end{axis}
    \end{tikzpicture}
    {\scriptsize Data space: three candidate lines.\par}
  \end{minipage}\hfill
  \begin{minipage}[t]{0.47\linewidth}
    \centering
    \begin{tikzpicture}
      \begin{axis}[mlaxis, width=\linewidth, height=4.4cm,
        xlabel={$w$}, ylabel={$J(w)$}, xmin=-0.6, xmax=2.6, ymin=-0.3, ymax=5.6]
        \addplot[MLBlue,very thick,domain=-0.6:2.6,samples=80]
          {14/6*(x-1)^2};
        \addplot[only marks,mark=*,mark size=2pt,MLRose,forget plot]
          coordinates {(0.5,0.5833)};
        \addplot[only marks,mark=*,mark size=2pt,MLBlue,forget plot]
          coordinates {(1,0)};
        \addplot[only marks,mark=*,mark size=2pt,MLAmber,forget plot]
          coordinates {(1.5,0.5833)};
        \node[font=\tiny,anchor=south east] at (axis cs:0.5,0.60) {$0.583$};
        \node[font=\tiny,anchor=south west] at (axis cs:1.5,0.60) {$0.583$};
        \node[font=\tiny,anchor=south west] at (axis cs:1.02,0.05) {$0$};
      \end{axis}
    \end{tikzpicture}
    {\scriptsize Parameter space: the cost of each.\par}
  \end{minipage}
  \caption{The same three models in both spaces. A line on the left is a point
  on the right. The curve on the right is not a model --- it is the cost of
  every possible model, and it is what the optimizer walks down.}
  \label{fig:parabola}
\end{figure}

\begin{center}
\small
\renewcommand{\arraystretch}{1.15}
\begin{tabular}{@{}lrrrrrrr@{}}
\toprule
$w$    & $-0.5$ & $0.0$ & $0.5$ & $1.0$ & $1.5$ & $2.0$ & $2.5$\\
$J(w)$ & $5.2500$ & $2.3333$ & $0.5833$ & $0.0000$ & $0.5833$ & $2.3333$ & $5.2500$\\
\bottomrule
\end{tabular}
\end{center}

Two observations to carry forward. First, $J$ is symmetric about its minimum
here, so being wrong by $0.5$ in either direction costs the same. Second,
$J(1)=0$ only because these three points happen to be collinear; on real data
the minimum of $J$ is positive, and \cref{sec:low-cost} explains why chasing
$J=0$ is a mistake rather than a goal.

\section{The Cost as a Function of Both Parameters}
\label{sec:cost-two-params}

Now let $b$ vary again, and return to the four-point set. Expanding the square
in the definition of $J$ and collecting terms gives a quadratic in two
variables:
\begin{align*}
  J(w,b) &= \frac{1}{2m}\sum_{i=1}^{m}\bigl(w\xii{i}+b-\yii{i}\bigr)^{2}\\
         &= \frac{1}{2m}\Bigl[
              w^{2}\textstyle\sum_i (\xii{i})^{2}
            + 2wb\sum_i \xii{i}
            + m b^{2}\\
         &\qquad\qquad
            - 2w\textstyle\sum_i \xii{i}\yii{i}
            - 2b\sum_i \yii{i}
            + \sum_i (\yii{i})^{2}
            \Bigr].
\end{align*}
For the four-point set the six sums are $\sum x^2 = 30$, $\sum x = 10$, $m=4$,
$\sum xy = 43$, $\sum y = 14$ and $\sum y^2 = 62$, so
\begin{equation}
  J(w,b) \;=\; \frac{30w^{2} + 20wb + 4b^{2} - 86w - 28b + 62}{8}.
  \label{eq:t4-cost}
\end{equation}

Every term is either a square or a product of the two parameters, and the
quadratic part $30w^2 + 20wb + 4b^2$ is positive for every $(w,b)\neq(0,0)$.
A quadratic of that kind is an \term{elliptic paraboloid}: a bowl.
\Cref{fig:bowl} draws it.

\begin{figure}[htbp]
  \centering
  \begin{tikzpicture}
    \begin{axis}[width=0.72\linewidth, height=5.6cm,
      view={40}{32},
      xlabel={$w$}, ylabel={$b$}, zlabel={$J(w,b)$},
      tick label style={font=\tiny}, label style={font=\small},
      colormap={ml}{color=(MLLight) color=(MLTeal) color=(MLBlue)},
      samples=26, domain=0.2:3.0, y domain=-3:2]
      \addplot3[surf,shader=interp,opacity=0.93]
        {(30*x^2 + 20*x*y + 4*y^2 - 86*x - 28*y + 62)/8};
    \end{axis}
  \end{tikzpicture}
  \caption{The cost surface \cref{eq:t4-cost} for the four-point set. Horizontal
  axes are the parameters; height is the cost. There is one lowest point, no
  local minima and no flat regions --- a shape called \emph{convex}, and the
  reason the optimizer of \cref{ch:ch03} cannot get stuck.}
  \label{fig:bowl}
\end{figure}

\subsection{Reading a contour map}
\label{sec:contours}

A three-dimensional surface is hard to read precisely and impossible to draw for
more than two parameters. The standard alternative is to look straight down at
it and draw \term{contours}: curves joining all points of equal height, exactly
as a topographic map joins points of equal elevation.

\begin{figure}[htbp]
  \centering
  \begin{minipage}[t]{0.48\linewidth}
    \centering
    \begin{tikzpicture}
      \begin{axis}[width=\linewidth, height=5.2cm,
        xlabel={$w$}, ylabel={$b$},
        tick label style={font=\scriptsize}, label style={font=\small},
        view={0}{90}, domain=0.2:3.0, y domain=-3:2,
        colormap={ml}{color=(MLBlue) color=(MLTeal) color=(MLAmber)}]
        \costcontour{}{0.1,0.25,0.6,1.2,2.2,3.6}
          {(30*x^2 + 20*x*y + 4*y^2 - 86*x - 28*y + 62)/8}
        \addplot3[only marks,mark=*,mark size=2.2pt,MLBlue]
          coordinates {(1.6,-0.5,0)};
        \addplot3[only marks,mark=square*,mark size=2pt,MLRose]
          coordinates {(1.0,1.0,0)};
        \addplot3[only marks,mark=triangle*,mark size=2.6pt,MLAmber]
          coordinates {(2.2,-2.0,0)};
        \node[font=\tiny,MLBlue,anchor=north west]  at (axis cs:1.62,-0.6,0) {A};
        \node[font=\tiny,MLRose,anchor=south west]  at (axis cs:1.04,1.0,0)  {B};
        \node[font=\tiny,MLAmber,anchor=north west] at (axis cs:2.24,-2.0,0) {C};
      \end{axis}
    \end{tikzpicture}
    {\scriptsize Parameter space, seen from above.\par}
  \end{minipage}\hfill
  \begin{minipage}[t]{0.48\linewidth}
    \centering
    \begin{tikzpicture}
      \begin{axis}[mlaxis, width=\linewidth, height=5.2cm,
        xlabel={$x$}, ylabel={$y$}, xmin=0, xmax=5, ymin=-1, ymax=8,
        legend style={font=\tiny, at={(0.02,0.98)}, anchor=north west}]
        \addplot[only marks,mark=*,mark size=1.8pt,MLTeal]
          coordinates {(1,1)(2,3)(3,4)(4,6)};
        \addlegendentry{data}
        \addplot[MLBlue,very thick,domain=0.2:4.6]{1.6*x-0.5};
        \addlegendentry{A: $J=0.025$}
        \addplot[MLRose,thick,domain=0.2:4.6]{1*x+1};
        \addlegendentry{B: $J=0.250$}
        \addplot[MLAmber,thick,dashed,domain=0.2:4.6]{2.2*x-2};
        \addlegendentry{C: $J=0.250$}
      \end{axis}
    \end{tikzpicture}
    {\scriptsize Data space: the three lines.\par}
  \end{minipage}
  \caption{The same information twice. Point~A is the minimum, $J = 0.025$.
  Points~B and~C sit on the \emph{same} contour and therefore have exactly equal
  cost, $J = 0.250$, despite being very different lines --- B too shallow, C too
  steep. The cost is a single number and cannot distinguish between the two ways
  of being wrong.}
  \label{fig:contours}
\end{figure}

Four things to take from \cref{fig:contours}, all of which recur throughout the
book:

\begin{enumerate}
  \item \textbf{The centre of the innermost closed curve is the minimum.} That
        is where training is trying to get to.
  \item \textbf{Points on one contour are equally good.} B and C both cost
        $0.250$. If your optimizer reports a cost, it is telling you which
        contour you are on and nothing about which direction you are off in.
  \item \textbf{Closely spaced contours mean a steep slope}, widely spaced ones a
        shallow slope. \Cref{ch:ch03}'s method moves perpendicular to the
        contours, so contour spacing is exactly what determines how fast it
        moves.
  \item \textbf{The contours are ellipses, not circles.} For the four-point set
        the ratio of the long axis to the short one is about $7.5$; for the
        apartment data of \cref{ch:ch01} it is about $41$. Elongation is caused
        by the features' scale, it slows optimization badly, and
        \cref{ch:ch05} shows how to remove it.
\end{enumerate}

\begin{pitfall}{Which space am I in?}
In \cref{fig:contours}, the left panel's horizontal axis is $w$ and the right
panel's is $x$. A point on the left is a whole line on the right. Students who
are stuck on cost-function questions are, nine times out of ten, reading a
parameter-space plot as though it were data space --- looking for the training
points in the contour plot, or expecting the fitted line to appear there. The
training data is not in the contour plot. It has been consumed by the definition
of $J$.
\end{pitfall}

\section{Fitting by Hand: The Closed Form}
\label{sec:closed-form}

For this model we can find the minimum exactly, without any iteration, because a
smooth function is stationary at its minimum: both partial derivatives vanish
there.

\begin{derivation}{The least-squares line}
Differentiate $J(w,b) = \frac{1}{2m}\sum_i (w\xii{i}+b-\yii{i})^2$ with respect
to each parameter in turn, using the chain rule:
\[
  \frac{\partial J}{\partial b}
    = \frac{1}{m}\sum_{i=1}^{m}\bigl(w\xii{i}+b-\yii{i}\bigr),
  \qquad
  \frac{\partial J}{\partial w}
    = \frac{1}{m}\sum_{i=1}^{m}\bigl(w\xii{i}+b-\yii{i}\bigr)\xii{i}.
\]
Setting the first to zero and writing $\bar x = \frac1m\sum_i \xii{i}$,
$\bar y = \frac1m\sum_i \yii{i}$:
\begin{equation}
  w\bar x + b - \bar y = 0
  \qquad\Longrightarrow\qquad
  b = \bar y - w\bar x .
  \label{eq:b-star}
\end{equation}
So the fitted line passes through the point $(\bar x, \bar y)$, the centre of
mass of the data, whatever $w$ turns out to be.

Setting the second to zero gives
$\sum_i (w\xii{i}+b)\xii{i} = \sum_i \yii{i}\xii{i}$, that is
$w\sum_i (\xii{i})^2 + b\sum_i \xii{i} = \sum_i \xii{i}\yii{i}$. Substituting
\cref{eq:b-star} for $b$ and using $\sum_i \xii{i} = m\bar x$:
\[
  \begin{aligned}
    w\sum_i (\xii{i})^{2} + (\bar y - w \bar x)\,m\bar x
      &= \sum_i \xii{i}\yii{i}\\[2pt]
    \Longrightarrow\quad
    w\Bigl[\sum_i (\xii{i})^{2} - m\bar x^{2}\Bigr]
      &= \sum_i \xii{i}\yii{i} - m\bar x\bar y ,
  \end{aligned}
\]
and both bracketed quantities can be rewritten as sums of deviations from the
mean, giving the form worth memorising:
\begin{equation}
  w^{\star} = \frac{\sum_{i=1}^{m}\bigl(\xii{i}-\bar x\bigr)
                        \bigl(\yii{i}-\bar y\bigr)}
                   {\sum_{i=1}^{m}\bigl(\xii{i}-\bar x\bigr)^{2}},
  \qquad
  b^{\star} = \bar y - w^{\star}\bar x .
  \label{eq:closed-form}
\end{equation}
The numerator is (up to a factor $m$) the covariance of $x$ and $y$, the
denominator the variance of $x$. The best slope is the covariance divided by the
variance --- which says, in one line, that the model tilts in proportion to how
strongly the two variables move together, tempered by how spread out the input
is.
\end{derivation}

\begin{workedexample}{The four-point set, fitted by hand}
$\bar x = (1+2+3+4)/4 = 2.5$ and $\bar y = (1+3+4+6)/4 = 3.5$. Tabulate the
deviations:

\begin{center}
\small
\renewcommand{\arraystretch}{1.15}
\begin{tabular}{@{}rrrrrr@{}}
\toprule
$\xii{i}$ & $\yii{i}$ & $\xii{i}-\bar x$ & $\yii{i}-\bar y$ &
$(\xii{i}-\bar x)(\yii{i}-\bar y)$ & $(\xii{i}-\bar x)^2$\\
\midrule
$1$ & $1$ & $-1.5$ & $-2.5$ & $3.75$ & $2.25$\\
$2$ & $3$ & $-0.5$ & $-0.5$ & $0.25$ & $0.25$\\
$3$ & $4$ & $ 0.5$ & $ 0.5$ & $0.25$ & $0.25$\\
$4$ & $6$ & $ 1.5$ & $ 2.5$ & $3.75$ & $2.25$\\
\midrule
    &     &        &        & $\mathbf{8.00}$ & $\mathbf{5.00}$\\
\bottomrule
\end{tabular}
\end{center}

Hence $w^\star = 8.00/5.00 = 1.6$ and $b^\star = 3.5 - 1.6(2.5) = -0.5$: the
fitted line is $f(x) = 1.6x - 0.5$.

Its residuals and cost:
\begin{center}
\small
\begin{tabular}{@{}rrrr@{}}
\toprule
$\xii{i}$ & $\yii{i}$ & $\yhat^{(i)}$ & $r^{(i)}$\\
\midrule
$1$ & $1$ & $1.1$ & $+0.1$\\
$2$ & $3$ & $2.7$ & $-0.3$\\
$3$ & $4$ & $4.3$ & $+0.3$\\
$4$ & $6$ & $5.9$ & $-0.1$\\
\bottomrule
\end{tabular}
\end{center}
$\sum_i (r^{(i)})^2 = 0.01+0.09+0.09+0.01 = 0.20$, so
$J = 0.20/(2\cdot 4) = 0.025$, matching point~A of \cref{fig:contours}.

\textbf{Two checks worth doing every time.} The residuals sum to
$0.1-0.3+0.3-0.1 = 0$, and weighted by $x$ they also sum to
$0.1 - 0.6 + 0.9 - 0.4 = 0$. These are not coincidences: they are precisely the
two equations $\partial J/\partial b = 0$ and $\partial J/\partial w = 0$ that
we solved. If your fitted residuals fail either check, the fit is wrong.
\end{workedexample}

Applying \cref{eq:closed-form} to the twelve apartments of \cref{ch:ch01}
returns $w^\star = 3.2984$ and $b^\star = 4.3308$ with $J = 2.4845$ --- the
numbers quoted there, now derived rather than asserted.

\subsection{Then why does the rest of the book exist?}

A fair question: we have solved the problem exactly, in closed form, in half a
page. Three reasons this is not the end of the story, each of which motivates a
later chapter.

\begin{itemize}
  \item \textbf{More features.} With $n$ features, \cref{eq:closed-form}
        generalises to a matrix equation requiring the inversion of an
        $n\times n$ matrix, at a cost of about $n^{3}$ operations. At $n =
        10{,}000$ features --- unremarkable for text or images --- that is out
        of reach. \Cref{ch:ch04} derives it anyway, and says exactly when to use
        it.
  \item \textbf{The matrix may not be invertible}, for instance when two
        features are duplicates of each other. Also \cref{ch:ch04}.
  \item \textbf{Most cost functions have no closed form at all.} The moment we
        change the model --- to logistic regression in \cref{ch:ch06}, or to a
        neural network in \cref{ch:ch09} --- setting the derivatives to zero
        produces equations no one can solve in closed form. An iterative method
        that needs only the derivatives, and never solves anything exactly, is
        the only approach that survives. That method is \cref{ch:ch03}'s
        gradient descent, and this chapter's bowl is the landscape it walks on.
\end{itemize}

\section{Least Squares as Projection}
\label{sec:projection}

The derivation above is correct but unilluminating; the geometry behind it is
worth five minutes because it makes the two residual checks obvious rather than
lucky.

Work in $\Real{m}$, with one dimension per training example. This space is not
the plane of the scatter plot. In it the observations form a single
vector $\vy = (\yii{1},\dots,\yii{m})$, the inputs form a vector
$\vx = (\xii{1},\dots,\xii{m})$, and $\mathbf{1} = (1,\dots,1)$. The model's
predictions are then
\[
  \hat{\vy} \;=\; w\,\vx + b\,\mathbf{1},
\]
which, as $w$ and $b$ range over all real numbers, sweeps out exactly the plane
through the origin spanned by $\vx$ and $\mathbf{1}$. Call it $\Hspace$: it is
the hypothesis space, drawn.

Minimising $J = \frac{1}{2m}\|\hat{\vy} - \vy\|^2$ therefore means finding the
point of that plane closest to $\vy$. Geometry answers this without any
calculus: the closest point is the \term{orthogonal projection} of $\vy$ onto
the plane, and the vector joining them --- the residual $\mathbf{r}$ --- is
perpendicular to the plane.

\begin{figure}[htbp]
  \centering
  \begin{tikzpicture}[scale=1.15, every node/.style={font=\small}]
    % the hypothesis plane, spanned by x and 1
    \fill[MLLight,draw=MLBlue!45]
      (-3.0,-0.9) -- (3.0,-1.7) -- (4.1,0.8) -- (-1.9,1.6) -- cycle;
    \node[MLBlue!70,anchor=north east,font=\footnotesize] at (4.0,0.7)
      {$\Hspace = \operatorname{span}\{\vx,\mathbf{1}\}$};
    \coordinate (O)   at (-1.2,-0.20);
    \coordinate (P)   at ( 1.70, 0.05);
    \coordinate (Y)   at ( 1.70, 2.60);
    \coordinate (X)   at ( 0.20,-0.95);
    \coordinate (ONE) at (-1.75, 0.85);
    \draw[-{Stealth},MLGray]      (O) -- (X)
      node[anchor=north,font=\footnotesize] {$\vx$};
    \draw[-{Stealth},MLGray]      (O) -- (ONE)
      node[anchor=east,font=\footnotesize] {$\mathbf{1}$};
    \draw[-{Stealth},MLBlue,very thick] (O) -- (Y)
      node[midway,anchor=south east,font=\footnotesize] {$\vy$};
    \draw[-{Stealth},MLTeal,very thick] (O) -- (P)
      node[midway,anchor=north,font=\footnotesize] {$\hat{\vy}=w\vx+b\mathbf{1}$};
    \draw[-{Stealth},MLRose,very thick] (Y) -- (P)
      node[midway,anchor=west,font=\footnotesize] {$\mathbf{r}=\hat{\vy}-\vy$};
    % right angle between the residual and the plane
    \draw[MLRose,thin] ($(P)+(0,0.30)$) -- ($(P)+(-0.30,0.33)$)
                       -- ($(P)+(-0.30,0.03)$);
    \fill[MLGray] (O) circle (1.2pt);
    \fill[MLTeal] (P) circle (1.4pt);
    \fill[MLBlue] (Y) circle (1.4pt);
  \end{tikzpicture}
  \caption{Least squares in $\Real{m}$. Every model the hypothesis space
  contains lies in the plane spanned by $\vx$ and $\mathbf{1}$. The observation
  vector $\vy$ generally does not. The best model is the point of the plane
  nearest $\vy$, and the residual is perpendicular to the plane.}
  \label{fig:projection}
\end{figure}

Perpendicular to the plane means perpendicular to everything in it, and in
particular to the two vectors that span it:
\[
  \mathbf{r}\cdot\mathbf{1} = 0
  \quad\text{i.e.}\quad \sum_i r^{(i)} = 0,
  \qquad\qquad
  \mathbf{r}\cdot\vx = 0
  \quad\text{i.e.}\quad \sum_i r^{(i)}\xii{i} = 0 .
\]
These are the two checks from the worked example, and they are the same two
equations obtained by setting the partial derivatives to zero. Calculus and
geometry give the identical answer, as they must.

\begin{keyinsight}{What fitting actually does}
Fitting does not find the ``true'' line. It finds the point of the hypothesis
space nearest the data, in the squared-distance sense. If the truth is not in
the hypothesis space --- if the data is curved and $\Hspace$ contains only
straight lines --- fitting will still return the nearest line, cheerfully and
without complaint. The distance that remains, $\|\mathbf{r}\|$, is the part of
the data your model cannot represent, and \cref{ch:ch10} gives it a name:
\emph{bias}.
\end{keyinsight}

\section{A Low Cost Is Not a Good Model}
\label{sec:low-cost}

The cost is the quantity we minimise. It is not the quantity we care about, and
the gap between the two is where most beginners' models go wrong.

\subsection{The cost has no scale}

$J = 2.4845$ was a good result in \cref{ch:ch01} and $J = 0.025$ is a good
result here, but neither number means anything on its own: $J$ carries the
squared units of $y$, so changing the units of the target changes $J$ by a
constant factor while changing nothing about the model. You cannot compare $J$
across datasets, and you cannot conclude anything from a single value of it. Use
a baseline, as \cref{ch:ch01} insisted, or report the root mean squared error
$\sqrt{2J}$, which at least has the same units as $y$.

\subsection{Squares make outliers loud}

Because the cost squares each residual, one badly wrong observation can dominate
the sum and drag the whole fit toward itself.

\begin{workedexample}{One bad point moves the line}
Add a single erroneous fifth observation to the four-point set, $(5, 30)$ ---
say a data-entry slip that put an extra digit on a price. Refitting with
\cref{eq:closed-form} gives

\begin{center}
\small
\begin{tabular}{@{}lrrrr@{}}
\toprule
& $w$ & $b$ & $\yhat$ at $x=1$ & $\yhat$ at $x=4$\\
\midrule
without the outlier & $1.6$ & $-0.5$ & $1.1$  & $5.9$\\
with the outlier    & $6.1$ & $-9.5$ & $-3.4$ & $14.9$\\
\midrule
observed $y$        &       &        & $1$    & $6$\\
\bottomrule
\end{tabular}
\end{center}

One point in five has nearly quadrupled the slope and made the model predict a
negative value where the data says $1$. The four original points are now fitted
far worse than before, and the fit is worse \emph{everywhere} rather than just
near the outlier. Nothing in the procedure noticed or objected: $J$ was
faithfully minimised throughout.
\end{workedexample}

\begin{pitfall}{Minimising $J$ is not the same as fitting well}
The optimizer answers exactly the question you asked: which member of
$\Hspace$ minimises this cost on this data. It cannot tell you that
\begin{itemize}
  \item an observation is a typo, or a different kind of thing from the rest;
  \item the relationship is curved, so no line will do (\cref{ch:ch05});
  \item the model fits the training data beautifully and new data terribly
        (\cref{ch:ch07});
  \item the input you fitted on will not be available at prediction time
        (\cref{ch:ch21}).
\end{itemize}
Always plot the data and the fitted model together, and always plot the
residuals against $x$. A residual plot showing structure --- a curve, a fan, a
cluster far from the rest --- is telling you that the hypothesis space is wrong,
and no amount of further optimization will fix that.
\end{pitfall}

\begin{labbox}{Lab 2 --- The cost function}
Implement \code{cost(x, y, w, b)} for the four-point set and confirm
$J(1.6,-0.5) = 0.0250$, $J(1,1) = 0.2500$ and $J(2.2,-2) = 0.2500$. Plot $J(w)$
with $b=0$ for the three-point set and check the vertex against the table in
\cref{sec:cost-one-param}. Then produce a contour plot matching
\cref{fig:contours}, and implement \cref{eq:closed-form} to recover
$w^\star = 1.6$, $b^\star = -0.5$. Finally, verify both residual identities to
machine precision: $\sum_i r^{(i)}$ and $\sum_i r^{(i)}\xii{i}$ should each come
out below $10^{-12}$ in magnitude, not merely small.
\end{labbox}

\section*{Chapter Summary}
\addcontentsline{toc}{section}{Chapter Summary}

\begin{itemize}
  \item The univariate model is $\fwb(x) = wx+b$; $w$ tilts the line and $b$
        raises it. Each pair $(w,b)$ is one point in parameter space and one
        line in data space, and keeping the two spaces distinct is essential.
  \item The residual is $r^{(i)} = \fwb(\xii{i}) - \yii{i}$, and the cost is
        $J(w,b) = \frac{1}{2m}\sum_i (r^{(i)})^2 = \frac{1}{2m}\|\mathbf{r}\|^2$.
  \item Squares are used because they are symmetric, differentiable, admit a
        closed-form solution, and correspond to Gaussian noise; the $1/m$ makes
        the cost independent of dataset size; the $\tfrac12$ cancels the $2$
        from differentiation and changes nothing else.
  \item With one parameter $J$ is an upward-opening parabola; with two it is a
        bowl. Both are convex, so there is exactly one minimum.
  \item Contours join parameter pairs of equal cost. The innermost contour
        surrounds the minimum, points on one contour are equally good, and
        elongated elliptical contours signal features on very different scales.
  \item The minimum is available in closed form:
        $w^\star = \sum_i (\xii{i}-\bar x)(\yii{i}-\bar y) / \sum_i (\xii{i}-\bar
        x)^2$ and $b^\star = \bar y - w^\star\bar x$. The fitted line always
        passes through $(\bar x, \bar y)$, and its residuals satisfy
        $\sum_i r^{(i)} = 0$ and $\sum_i r^{(i)}\xii{i} = 0$.
  \item Geometrically, fitting is the orthogonal projection of $\vy$ onto the
        plane spanned by $\vx$ and $\mathbf{1}$; the two residual identities are
        that perpendicularity.
  \item A small $J$ is not by itself evidence of a good model: $J$ has no
        absolute scale, squaring amplifies outliers, and the cost is blind to
        every question except the one it was asked.
\end{itemize}

\section*{Exercises}
\addcontentsline{toc}{section}{Exercises}

\exercisesection{Short Questions}

\sq State the univariate linear regression model and say, in one sentence each,
what $w$ and $b$ control.

\sq Define the residual of a training example. Under what circumstances is it
positive?

\sq Write down the squared-error cost function and state the units of $J$ when
$y$ is measured in kilograms.

\sq Give two reasons for squaring the residuals rather than taking their
absolute value.

\sq Why does the $\tfrac12$ in $J$ not change which line is chosen?

\sq A dataset is duplicated so that every example appears twice. What happens to
$\sum_i (r^{(i)})^2$, and what happens to $J$?

\sq Sketch $J(w)$ for a model with $b$ fixed. Why can this curve have no local
minimum other than the global one?

\sq Two parameter pairs lie on the same contour of $J$. What does that tell you
about them, and what does it not tell you?

\sq The contours of $J$ for some dataset are long thin ellipses. What does this
say about the data, and which later chapter addresses it?

\sq State the two identities satisfied by the residuals of a least-squares fit,
and say what they mean geometrically.

\exercisesection{Analytical Problems}

\ap For the four-point set of \cref{sec:notation}, compute $J(2,0)$ and
$J(1,0.5)$ by hand, showing the residual for each example. Which of the two
lines is better, and by how much?

\ap Using \cref{eq:closed-form}, fit a line to the three points $(0,1)$,
$(1,2)$, $(2,6)$. Report $w^\star$, $b^\star$ and $J$, and verify both residual
identities.

\ap Show that the value of $b$ minimising $J$ for a fixed $w$ is
$b = \bar y - w\bar x$, and hence that every least-squares line passes through
$(\bar x, \bar y)$, regardless of the data.

\ap Suppose every target is multiplied by a constant $c > 0$, so that
$\yii{i} \mapsto c\,\yii{i}$. Show that the optimal parameters become
$(cw^\star, cb^\star)$ and that the minimum cost becomes $c^2$ times its
original value. What does this imply about comparing $J$ between two models
whose targets are measured in different units?

\ap Expand $J(w,b)$ for a general dataset into the form
$Aw^2 + Bwb + Cb^2 + Dw + Eb + F$ and express each of the six coefficients in
terms of $m$ and the sums $\sum x$, $\sum x^2$, $\sum y$, $\sum y^2$,
$\sum xy$. Verify your expressions against \cref{eq:t4-cost}.

\ap Consider fitting the constant model $f(x) = b$ (that is, $w$ forced to zero)
to the four-point set. Find the optimal $b$ and the resulting $J$, and compare
with the $J = 0.025$ of the full model. What fraction of the cost did the slope
remove?

\exercisesection{Design Problems}

\dpr A colleague proposes replacing the squared-error cost with
$J' = \frac{1}{m}\sum_i |r^{(i)}|$ for a dataset known to contain occasional
data-entry errors. Set out the argument for and against, covering the effect on
outliers, on the existence of a closed-form solution, and on the optimizer of
\cref{ch:ch03}. What would you recommend, and what would you check first?

\dpr You are fitting a model that predicts a patient's dose of a drug from their
body mass. Over-prediction and under-prediction are not equally harmful. Propose
a modification to the cost function that reflects this, state what property of
$J$ your modification preserves or destroys, and describe how you would choose
its parameter.

\dpr A model is to be reported to a non-technical audience who will ask ``how
accurate is it?''. You have $J$, the root mean squared error, and the mean
absolute error available. Choose one, justify the choice, and write the single
sentence you would put in the report.

\dpr Design a small dataset --- at most six points --- for which the
least-squares line is a visibly poor description of the data even though it is
the exact minimiser of $J$. Explain what feature of your dataset causes this,
and what a residual plot would show.

\exercisesection{Programming Exercises}

\pe Implement \code{cost(x, y, w, b)} without a loop over examples, and
reproduce the table in \cref{sec:cost-one-param} and the three values marked in
\cref{fig:contours}.

\pe Plot $J(w)$ for the four-point set with $b$ fixed at $-0.5$, over
$w \in [0, 3]$. Mark the minimum. Then plot it again with $b$ fixed at $0$ and
explain, in one sentence, why the two minima occur at different values of $w$.

\pe Draw a contour plot of $J(w,b)$ for the four-point set by evaluating $J$ on
a grid and calling your plotting library's contour routine. Overlay the three
points A, B and C of \cref{fig:contours}. Then repeat for the twelve apartments
of \cref{ch:ch01} and comment on the difference in shape.

\pe Implement \cref{eq:closed-form} and verify it on both the four-point set and
the apartment data. Then write a brute-force alternative that evaluates $J$ over
a fine grid of $(w,b)$ and reports the best pair found. How fine must the grid be
to agree with the closed form to three decimal places, and how many evaluations
is that?

\exercisesection{Self-Test}

\qz $J(w,b)$ is a function of: (a)~the training data; (b)~the parameters;
(c)~the predictions; (d)~the residuals only.

\qz Doubling every residual multiplies $J$ by: (a)~$2$; (b)~$4$;
(c)~$\sqrt{2}$; (d)~it depends on $m$.

\qz The $\tfrac12$ in the cost function exists to: (a)~halve the error;
(b)~cancel the $2$ from differentiating the square; (c)~keep $J$ below $1$;
(d)~average over both parameters.

\qz Points on the same contour of $J$ have equal: (a)~slope; (b)~intercept;
(c)~cost; (d)~residuals.

\qz The least-squares line always passes through: (a)~the origin; (b)~the first
training example; (c)~the point $(\bar x, \bar y)$; (d)~none of these.

\qz If $\sum_i r^{(i)} \neq 0$ for a supposedly least-squares fit, then:
(a)~the data is noisy; (b)~the fit is wrong; (c)~$m$ is too small; (d)~the cost
needs a $\tfrac12$.

\qz The contours of $J$ are circles rather than ellipses when: (a)~the data is
noise-free; (b)~$m$ is large; (c)~the feature is centred and suitably scaled;
(d)~never.

\qz A model achieves $J = 0$ on its training set. This means: (a)~the model is
correct; (b)~the model will generalise; (c)~every training point lies exactly on
the fitted line; (d)~the cost function is misconfigured.

\furtherreading

\textcite{james2021} §3.1 covers simple linear regression, including the
closed-form coefficients derived in \cref{sec:closed-form}; read §3.1.1 and
§3.1.2 now and return to §3.1.3 after \cref{ch:ch10}, since it concerns
assessing a fit rather than obtaining one. \textcite{bishop2006} §1.1 sets up
the same cost through polynomial curve fitting and is worth reading for the
alternative framing, while §3.1.1 gives the geometric projection picture of
\cref{sec:projection} in full generality.

\bigskip
\noindent\emph{Next:} \Cref{ch:ch03} keeps the bowl and throws away the closed
form. We ask how a program that can measure the local slope, and nothing else,
can find the bottom of a surface it cannot see.
